\section{3D Heterogeneous Integration \& TSVs}

\subsection{2D vs 2.5D vs 3D Integration}
Traditional 2D integration puts transistors and functional parts on one flat silicon chip. This method works well. 
Is very old but as data moves between different parts it can slow things down and use more power. 
In 2.5D integration several chips sit next to each other on a base called an interposer. The interposer 
connects them closely without stacking the chips on top of each other. In integration the chips or layers 
are stacked up like a tower. Connections between the layers happen through paths made with technologies 
such, as Through-Silicon Vias (TSVs) [6] [7].

\subsection{Heterogeneous Integration}
Heterogeneous integration brings together components made using technologies, materials, process nodes or 
functions into a single system. For example, a package might combine processor logic, memory, sensors, 
radio-frequency parts, or accelerators that were produced using manufacturing processes [8]. This method 
supports the "More-than-Moore" idea, where system performance and functionality improve not by shrinking 
transistors but through smarter integration [9]. Chiplet-based architectures are an example of this since 
each chiplet can be designed and optimized separately before being assembled into a larger system.

\subsection{Through-Silicon Vias}
A Through-Silicon Via or TSV is an electrical connection that goes through a silicon die or wafer. 
These vias are usually filled with a material, most often copper and are used to link circuits or 
stacked dies in a vertical direction. TSVs allow for shorter communication paths compared to traditional 
package-level connections. The technology involves key steps like wafer thinning creating the via holes 
adding insulation depositing metal layers and bonding. Each of these steps needs to be controlled. 
Any defects in the TSVs can lead to problems, with performance and can reduce the long-term reliability 
of the device.

\subsection{Performance Advantages}
3D integration can greatly boost how many connections fit in an area while making the space between 
processing parts and memory parts much smaller. When connections are shorter data moves faster. 
Uses less energy. Stacking components vertically means more features can fit into a space [6]. 
High-Bandwidth Memory or HBM shows this idea in action. It uses memory chips stacked on top of 
each other with vertical wires called TSVs. These allow data transfer between the memory and the 
processor or logic chip [10]. This kind of setup also lets different parts be built using the 
manufacturing methods available, for each component.

\subsection{Challenges}
Despite its benefits 3D integration brings a lot of problems. Getting rid of heat is harder because the dies that 
are stacked inside can stop heat from moving to the outside of the package. Thermal stress, bending, issues with 
through silicon vias, production checking for faults and getting enough good parts are also problems. These are 
mentioned in sources [8] and [9]. Also the steps needed to make the wafers thin stick them together. Create the 
through silicon vias can cost more money. So making 3D integration work well needs planning of the design the case, 
around the chip handling the heat making the parts and making sure they last a long time.
